
\def\PUDcodigo{EME111}

\def\PUDcodacad{04506.28}

\def\PUDdisciplina{Resistência dos Materiais II}

\def\PUDsemestre{S6}

\def\PUDhoras{80}

%NOVOS ---------------------------------------------------
\def\PUDhorasTeoria{80}
\def\PUDhorasPratica{0}

\def\PUDinterECA{
\hyperref[PUD:EME109]{Resistência dos Materiais I (04506.26)}

\hyperref[PUD:EME119]{Sistemas Mecânicos (04506.35)}
}

\def\PUDatuaisECA{---}

\def\PUDrecursos{Projetor multimídia; Quadro branco e pincel.}

%----------------------------------------------------------

\def\PUDcreditos{4}

\def\PUDprerequisito{ \hyperref[PUD:EME109]{EME109} }

\def\PUDpreacad{ \hyperref[PUD:EME109]{Resistência dos Materiais I (04506.26)} }

\def\PUDobjetivos{Conhecer campos de deslocamentos em problemas hiperestáticos através de diversos métodos;  Conhecer os conceitos de grau de liberdade, discretização, matrizes estruturais, condições de contorno, nós e elementos e operações de análise estrutural matricial através do método de elementos finitos;  Conhecer as teorias para alguns modos de falha: flambagem, plastificação em flexão e resistência a fadiga de metais.}

\def\PUDementa{Cargas Combinadas;  Transformação de Tensão;  Transformação da Deformação;  Deflexão em Vigas e Eixos;  Flambagem de Colunas;  Métodos de Energia.}

\def\PUDconteudo{ & Cargas combinadas: Vasos de pressão de paredes finas, estados de tensão causados por tensões combinadas. 
\\ 
 & Transformação de tensão: Transformação de tensões no plano, equações gerais de transformações de tensão no plano, tensões principais e tensão de cisalhamento máxima, Círculo de Mohr – tensão no plano, tensões em eixos causadas por cargas axiais e torcionais, variações de tensão ao longo de uma viga, tensão de cisalhamento máxima absoluta.  
\\ 
 & Transformação da deformação: deformação plana, Círculo de mohr – plano de deformações, deformação por cisalhamento máximo absoluto, rosetas de deformação, relação entre materiais e suas propriedades. 
\\ 
 & Deflexão em vigas e eixos: A linha elástica, inclinação e deslocamento, funções de descontinuidade, método da superposição, vigas e eixos estaticamente indeterminados
\\ 
 & Flambagem de colunas: carga crítica, colunas com vários tipos de apoio, a fórmula da secante, flambagem inelástica, projeto de colunas. 
\\ 
 & Métodos de energia: Trabalho externo e energia de deformação, conservação de energia, carga de impacto, princípio do trabalho virtual, método das forças virtuais, Teorema de Castigliano. }

\def\PUDmetodo{Aulas expositórias que podem ser teóricas e/ou práticas, onde as práticas no laboratório serão marcadas no decorrer da disciplina.}

\def\PUDavalia{A avaliação será feita com aplicação de provas teóricas e/ou práticas, além da possibilidade de inclusão de trabalhos e seminários no decorrer da disciplina.}

\def\PUDbibbasica{ & \href{www.biblioteca.ifce.edu.br/index.asp?codigo_sophia=24229}{Hibbeler}, R. C.  Resistência dos Materiais.  5ª ed. São Paulo, SP: Prentice Hall, 2004.  670p.  ISBN: 9788587918673. 
\\ 
& \href{biblioteca.ifce.edu.br/index.asp?codigo_sophia=24321}{Beer}, Ferdinand P. Resistência dos materiais. 3. ed. São Paulo: Pearson Makron Books, 2008. 1255 p. Inclui índice. ISBN 9788534603447.
%\href{www.biblioteca.ifce.edu.br/index.asp?codigo_sophia=24421}{Callister} Jr., W. D.; Rethwisch, D. G.  Ciência e Engenharia de Materiais - Uma Introdução.  8ª ed.  Rio de Janeiro, RJ: LTC, 2012.  817p.  ISBN: 9788521621249.
\\ 
& \href{www.biblioteca.ifce.edu.br/index.asp?codigo_sophia=61931}{Johnston Jr.}, E. R.;  Beer, F. P.;  Mecânica dos Materiais.  7ª ed.  Porto Alegre, RS: AMGH, 2015.  838p.  ISBN: 9788580554984. 
   %\href{www.biblioteca.ifce.edu.br/index.asp?codigo_sophia=61931}{Johnston} Jr., E. R.; Beer, F. P.  Mecânica dos Materiais.  7ª ed.  Porto Alegre, RS: AMGH, 2015.  838p.  ISBN: 9788580554984. 
   }

\def\PUDbibcomplementar{ & \href{www.biblioteca.ifce.edu.br/index.asp?codigo_sophia=59395}{Pereira}, Celso Pinto Morais.  Mecânica dos materiais avançada.  E-book.  Interciência.  434p.  ISBN: 9788571933347.
\\
& \href{biblioteca.ifce.edu.br/index.asp?codigo_sophia=40282}{Norton}, Robert L. Projeto de máquinas: uma abordagem integrada. 4. ed. Porto Alegre: Bookman, 2013. 1028 p. ISBN 9788582600221.
\\
& \href{biblioteca.ifce.edu.br/index.asp?codigo_sophia=63128}{Juvinall}, Robert C.; Marshek, Kurt M. Fundamentos do projeto de componentes de máquinas. 5. ed. Rio de Janeiro: LTC, 2016. 562 p., il. ISBN 9788521630098.
\\
& \href{biblioteca.ifce.edu.br/index.asp?codigo_sophia=23902}{Collins}, Jack A. Projeto mecânico de elementos de máquinas: uma perspectiva de prevenção da falha. Rio de Janeiro: LTC, 2006. 740 p. Inclui bibliografia e indice. ISBN 8521614756
\\
& \href{www.biblioteca.ifce.edu.br/index.asp?codigo_sophia=42347}{Shackelford}, James F.  Ciências dos materiais.  6 ed.  São Paulo, SP: Pearson, 2012.  556p.  ISBN: 9788576051602.
%&  BEER, F.  P.;  DEWOLF, J.  T.;  Resistência dos Materiais.  4ª edição.  São Paulo.  Editora  MCGRAW HILL - ARTMED.  2006.  774p.  ISBN 9788586804830. 
%\\ 
%  & \href{www.biblioteca.ifce.edu.br/index.asp?codigo_sophia=44596}{Garcia}, Amauri.  Ensaios dos materiais.  2 ed.  Rio de Janeiro, RJ: LTC, 2012.  365p.  ISBN: 9788521620679.
%\\ 
 %& \href{www.biblioteca.ifce.edu.br/index.asp?codigo_sophia=62306}{Nunes}, Laerce de Paula.  Fundamentos de resistência à corrosão.  Rio de Janeiro, RJ: Interciência, 2007.  330p.  ISBN: 9788571931626.
%\\
% & \href{www.biblioteca.ifce.edu.br/index.asp?codigo_sophia=42065}{Bauer}, L. A. F.  Materiais de construção.  v.1.  5º ed.  Rio de Janeiro, RJ: LTC, 2012.  471p.  ISBN: 9788521612490.
%\\
 %&  SILVA, L.  F.  M.;  GOMES, J.  F.  S.  Introdução à Resistência dos Materiais. 1ª edição, Portugal.  Editora Publindustria.  2010.  308p  ISBN 9789728953553. 
 }

\def\PUDautor{Venceslau Xavier de Lima Filho}

\def\PUDdata{2013-06-16}

\def\PUDrevisao{2}

\def\PUDrevisor{Venceslau Xavier de Lima Filho}

\def\PUDdatarevisao{2019-05-15}

\PUDcorpo
